\subsection*{The Standard One-Based System}

\begin{topicbox}{The Usual One-Based Carrier}
\begin{exposition}
The familiar natural-number model is the reference example, but it is only one
realization of the Peano axioms. The carrier is the usual one-based set of
natural numbers, the distinguished element is $1$, and successor is the usual
next-number operation.
\end{exposition}

\begin{example*}[The Usual Natural Numbers]
Let
\[
P=\mathbb{N}, \qquad e=1,
\]
and define
\[
S:P\to P,
\qquad
S(n)=n+1.
\]
Then $(P,S,e)=(\mathbb{N},S,1)$ is the usual one-based Peano system.
\end{example*}

\begin{remark*}[Structure]
The Peano data are
\[
\bigl(P,S,e\bigr)
=
\bigl(\mathbb{N},\,n\mapsto n+1,\,1\bigr).
\]
\end{remark*}

\begin{remark*}[Interpretation]
This is the familiar model, but it is not the only possible carrier for a
Peano system. The remaining examples use different elements while preserving
the same abstract successor pattern.
\end{remark*}
\end{topicbox}

